\documentclass[12pt]{scrartcl}

\usepackage{amsmath,amssymb}
\usepackage{fullpage}
\usepackage{enumitem}
\usepackage{hyperref}
\usepackage{listings}
\usepackage{xcolor}

\setlength{\parindent}{0pt}

\lstset{
    language=Python,
    basicstyle=\ttfamily\small,
    keywordstyle=\color{blue},
    commentstyle=\color{gray},
    stringstyle=\color{teal},
    frame=single,
    columns=fullflexible,
    keepspaces=true,
    showstringspaces=false
}

\begin{document}

\begin{center}
	\hrule
	\vspace{0.4cm}
	{\textbf{\large INF1008 Data Structures and Algorithms}}\\[0.2cm]
	Tutorial 1
\end{center}

\textbf{Name:} Timothy Chia\hspace{\fill} \textbf{Student ID:} 2501530\\

\hrule

\vspace{0.3cm}

\begin{enumerate}[label=\textbf{\arabic*.}]

	\item
	      Write an algorithm (pseudocode is sufficient) that returns the index of the first occurrence of the maximum element in an array \( s[0], \ldots, s[n-1] \).
	      Use only basic \texttt{for} loops and \texttt{if} statements.

	      \vspace{0.3cm}
	      \begin{lstlisting}
arr: list[int] = [3, 2, 1, 5, 3, 2, 1, 5]

def find_largest_index(a: list[int]) -> int:
  max_idx: int = 0 # start with first element as largest.

  # iterate remainder of array, O(n + 1)
  for i in range(1, len(a)):
    if a[i] > a[max_idx]:
      max_idx = i

  return max_idx

idx: int = find_largest_index(arr)

print(f"idx: {idx}")  # idx: 3
print(f"arr[{idx}]: {arr[idx]}")  # arr[3]: 5
\end{lstlisting}

	      %------------------------------------------------------------
	      \vspace{0.3cm}
	\item
	      Write an algorithm (pseudocode is sufficient) that reverses the order of elements in an array \( s[0], \ldots, s[n-1] \).
	      Use only basic \texttt{for} loops and \texttt{if} statements.

	      \vspace{0.3cm}
	      \begin{lstlisting}
# question 2
arr: list[int] = [10, 9, 8, 7, 6, 5, 4, 3, 2, 1]

def reverse(a: list[int]) -> list[int]:
  n: int = len(a) // 2

  for i in range(0, n):
    left = i
    right = (len(a) - i) - 1
    
    a[left], a[right] = a[right], a[left]

  return a

print(reverse(arr)) # [1, 2, 3, 4, 5, 6, 7, 8, 9, 10]
\end{lstlisting}

	      %------------------------------------------------------------
	      \newpage
	\item
	      Write an algorithm (pseudocode is sufficient) that outputs the largest and second-largest values in an array \( s[0], \ldots, s[n-1] \).
	      Assume \( n > 1 \) and that all elements are distinct.
	      Use only basic \texttt{for} loops and \texttt{if} statements.

	      \vspace{0.3cm}
	      \begin{lstlisting}
arr: list[int] = [8, 6, 4, 2, 9, 7, 5, 3, 1]

def find_largest_two(a: list[int]) -> tuple[int, int]:
  first_max_idx: int = 0
  second_max_idx: int = 1

  # edge case
  if a[second_max_idx] > a[first_max_idx]:
    first_max_idx, second_max_idx = second_max_idx, first_max_idx

  for i in range(2, len(a)):
    is_gt_first_max: bool = a[i] > a[first_max_idx]
    is_gt_second_max: bool = a[i] > a[second_max_idx]

    if is_gt_first_max:
      first_max_idx, second_max_idx = i, first_max_idx
    elif is_gt_second_max:
      second_max_idx = i
    else:
      continue

  return a[first_max_idx], a[second_max_idx]

first_largest, second_largest = find_largest_two(arr)

print(f"first_largest: {first_largest}")  # first_largest: 9
print(f"second_largest: {second_largest}")  # second_largest: 8
\end{lstlisting}

	      %------------------------------------------------------------
	      \newpage
	\item
	      Given a sorted array \( s[0], \ldots, s[n-1] \) where \( n > 1 \) and \( s[i] \leq s[i+1] \) for all valid \( i \),
	      write an algorithm (pseudocode is sufficient) that inserts a value \( x \) such that the array remains sorted.
	      Use only basic \texttt{for} loops and \texttt{if} statements.

	      \vspace{0.3cm}
	      \begin{lstlisting}
arr: list[int] = [1, 3, 5, 7, 9, 11, 13, 15]

def insert(a: list[int], x: int) -> list[int]:
  n: int = len(a)
  
  # iterate list and check if element to insert is 
  # lt current element then insert at position
  for i in range(0, n):
    if x < a[i]:
      a.insert(i, x)
      return a
      
  # element to insert is larger than all current elements  
  a.append(x)
  return a

insert(arr, 17)

print(arr)  # [1, 3, 5, 7, 9, 11, 13, 15, 17]
\end{lstlisting}

	      %------------------------------------------------------------
	      \newpage
	\item
	      Consider the following algorithm for finding the maximum element in an array:

	      \begin{lstlisting}
function arrayMax(a, n):
    currentMax = a[0]
    for i = 1 to n
        if a[i] > currentMax
            currentMax = a[i]
    return currentMax
\end{lstlisting}

	      Determine how many times the assignment statement
	      \texttt{currentMax = a[i]} executes in the best case and in the worst case.

	      \vspace{0.3cm}
	      \begin{lstlisting}
# question 5
worst_case: list[int] = [1, 2, 3, 4, 5]
best_case: list[int] = [5, 4, 3, 2, 1]

def max(a: list[int]) -> tuple[int, int]:
  max_idx = 0
  count = 0

  for i in range(1, len(a)):
    if a[i] > a[max_idx]:
      max_idx = i
      count += 1

  return a[max_idx], count

_, count = max(worst_case)
print(f"In the worst case, max_idx = i executed {count} times.")
# In the worst case, max_idx = i executed 4 (n - 1) times.

_, count = max(best_case)
print(f"In the best case, max_idx = i executed {count} times.")
# In the best case, max_idx = i executed 0 times.
\end{lstlisting}
\end{enumerate}
\end{document}
